\documentclass[11pt,a4paper]{ctexart}
\usepackage[margin=1in]{geometry}
\usepackage{booktabs}
\usepackage{longtable}
\usepackage{array}
\usepackage{tabularx}
\usepackage{hyperref}
\usepackage{enumitem}
\usepackage{float}
\usepackage{xcolor}

\hypersetup{colorlinks=true, linkcolor=blue, urlcolor=blue}
\setlength{\parskip}{0.4em}
\setlength{\parindent}{2em}

\title{MSWC English 数据集规模与 8:1:1 划分建议\\\large VE216 Topic 2 语音项目数据说明}
\author{基于本地原始数据与当前子集的统计结果}
\date{2026年4月25日}

\begin{document}
\maketitle

\section{文档目的}

这份说明文档回答四个问题：
\begin{enumerate}[nosep]
    \item MSWC English 原始数据到底有多大；
    \item 当前项目里的 160 词子集为什么会让人感觉“每个辅音的词数不够”；
    \item 如果明确要求使用 \textbf{8:1:1} 的 train/dev/test 比例，怎样分配更合理；
    \item 在当前课程项目规模下，究竟应该取多少词、每词取多少条音频。
\end{enumerate}

本文所有数字都来自本地现有文件：
\begin{itemize}[nosep]
    \item \path{216/project/topic2_speech/data/raw/mswc_en/metadata.json.gz}
    \item \path{216/project/topic2_speech/data/raw/mswc_en/en_splits.csv}
    \item \path{216/project/topic2_speech/data/raw/mswc_en/audio_en.tar.gz}
    \item \path{216/project/topic2_speech/data/raw/mswc_en_subset/selection_summary.json}
    \item \path{216/project/topic2_speech/data/raw/mswc_en_subset_600/selection_summary.json}
    \item \path{216/project/topic2_speech/src/prepare_mswc_subset.py}
\end{itemize}

\paragraph{2026-04-26 更新。}
本文最初是为“为什么要从 160 词迁移到 600 词”写的数据设计说明。现在这一步已经落地完成：脚本默认值已切到 600 / 40 / 5 / 5，\texttt{seed\_words\_per\_label} 也已经暴露成命令行参数，实际生成的子集位于 \path{data/raw/mswc_en_subset_600/}。因此，本文现在更适合作为“迁移理由 + 规模设计依据”来阅读，而不是未来时的提案。

\section{MSWC English 原始数据有多大}

\subsection{整体规模}

表~\ref{tab:raw-overview} 总结了当前本地 \textbf{MSWC English} 的规模。

\begin{table}[H]
    \centering
    \caption{MSWC English 原始数据整体规模}
    \label{tab:raw-overview}
    \begin{tabular}{>{\raggedright\arraybackslash}p{8.0cm}r}
        \toprule
        项目 & 数值 \\
        \midrule
        metadata 中记录的英文词种数 & 38,174 \\
        metadata 中记录的总音频数 & 6,624,343 \\
        \texttt{VALID=True} 的有效音频总数 & 6,621,345 \\
        有效 split 覆盖的英文词种数 & 38,150 \\
        train 有效样本数 & 5,264,225 \\
        dev 有效样本数 & 678,490 \\
        test 有效样本数 & 678,630 \\
        \path{metadata.json.gz} 文件大小 & 107.75 MB \\
        \path{en_splits.csv} 文件大小 & 1.24 GB \\
        \path{audio_en.tar.gz} 压缩包大小 & 34.85 GB \\
        本地原始目录总大小 & 36 GB \\
        \bottomrule
    \end{tabular}
\end{table}

这个规模说明：\textbf{MSWC English 作为原始数据源完全够大}。当前项目里“某些辅音只有少量单词可用”的问题，不是数据源不够，而是我们后续子集筛选太保守。

\subsection{词频分布是明显的长尾分布}

MSWC English 不是“每个单词都差不多多”的均匀词表，而是很典型的长尾分布。表~\ref{tab:quantiles} 给出词频分位数。

\begin{table}[H]
    \centering
    \caption{MSWC English 英文词频分位数}
    \label{tab:quantiles}
    \begin{tabular}{cccccccc}
        \toprule
        10\% & 25\% & 50\% & 75\% & 90\% & 95\% & 99\% & 最大值 \\
        \midrule
        5 & 7 & 17 & 68 & 238 & 517 & 2350 & 140,555 \\
        \bottomrule
    \end{tabular}
\end{table}

这组数字的意思是：
\begin{itemize}[nosep]
    \item 一半的英文单词，在 metadata 里最多只有 17 条音频；
    \item 只有头部少量高频词，才能达到几百、几千甚至十几万条；
    \item 所以如果我们把每词最低门槛设得很高，就会快速丢掉大量中尾部词种。
\end{itemize}

表~\ref{tab:topwords} 给出最常见的 10 个英文词。

\begin{table}[H]
    \centering
    \caption{MSWC English 中最常见的 10 个英文词}
    \label{tab:topwords}
    \begin{tabular}{lr}
        \toprule
        单词 & 音频数 \\
        \midrule
        was & 140,555 \\
        and & 116,658 \\
        you & 87,515 \\
        that & 76,472 \\
        are & 62,189 \\
        his & 59,584 \\
        this & 55,835 \\
        for & 45,417 \\
        have & 42,396 \\
        what & 37,490 \\
        \bottomrule
    \end{tabular}
\end{table}

因此，设计课程项目子集时，真正要平衡的是两件事：
\begin{enumerate}[nosep]
    \item 每个词内部需要保留多少录音；
    \item 我们想覆盖多少个不同的词种。
\end{enumerate}

对本题来说，第二件事更关键，因为我们研究的是“某个浊辅音在不同单词中是否有稳定频谱模式”，而不是“同一个单词重复多少次”。

\section{当前 160 词子集的问题在哪里}

当前项目现成子集来自 \path{src/prepare_mswc_subset.py} 的默认配置：
\begin{itemize}[nosep]
    \item \texttt{target-words = 160}
    \item 每词 \texttt{train/dev/test = 60/8/8}
    \item 总文件数 \texttt{160 * (60+8+8) = 12160}
\end{itemize}

其当前统计见表~\ref{tab:current-subset}。

\begin{table}[H]
    \centering
    \caption{当前 160 词子集的规模}
    \label{tab:current-subset}
    \begin{tabular}{>{\raggedright\arraybackslash}p{7.5cm}r}
        \toprule
        项目 & 数值 \\
        \midrule
        词数 & 160 \\
        总 WAV 数 & 12,160 \\
        当前每词配额 & 60 / 8 / 8 \\
        train / dev / test 总数 & 9,600 / 1,280 / 1,280 \\
        子集目录大小 & 390 MB \\
        平均单个 WAV 大小 & 32,072 Bytes $\approx$ 31.3 KB \\
        可用浊辅音标签数 & 15 \\
        \bottomrule
    \end{tabular}
\end{table}

这个配置有两个明显问题：
\begin{enumerate}[nosep]
    \item \textbf{比例不是严格的 8:1:1。} 60/8/8 更接近 7.5:1:1。
    \item \textbf{真正的瓶颈是词种数，而不是录音条数。} 每个词拿了 76 条录音，但一共只选 160 个词，所以很多辅音只覆盖到很少的不同单词。
\end{enumerate}

当前子集里，各标签对应的词数如表~\ref{tab:current-labels}。

\begin{table}[H]
    \centering
    \caption{当前 160 词子集的标签词数分布}
    \label{tab:current-labels}
    \small
    \begin{tabular}{cccccc}
        \toprule
        标签 & 词数 & 标签 & 词数 & 标签 & 词数 \\
        \midrule
        b & 16 & l & 32 & v & 16 \\
        d & 27 & m & 28 & w & 25 \\
        dh & 14 & n & 45 & y & 10 \\
        g & 11 & ng & 15 & z & 15 \\
        jh & 6 & r & 29 & zh & 4 \\
        \bottomrule
    \end{tabular}
\end{table}

这张表直接解释了“为什么感觉词不够用”：
\begin{itemize}[nosep]
    \item \texttt{/zh/} 只有 4 个词；
    \item \texttt{/jh/} 只有 6 个词；
    \item \texttt{/y/} 只有 10 个词；
    \item \texttt{/g/} 只有 11 个词。
\end{itemize}

这时模型很容易学成“这些单词长什么样”，而不是“这个辅音本身长什么样”。

\subsection{当前门槛本身就会压缩候选池}

当前脚本里，\texttt{train/dev/test} 既是最终抽样配额，也是候选词的最低门槛。也就是说，一个词只有在 train 至少 60、dev 至少 8、test 至少 8 时，才有资格进入候选池。

在本地重新统计后：
\begin{itemize}[nosep]
    \item 当前 \texttt{60/8/8} 门槛下，满足条件的候选词共有 \textbf{8,227} 个；
    \item 如果改成严格 8:1:1 且更温和的 \texttt{40/5/5}，候选词会增加到 \textbf{10,417} 个。
\end{itemize}

所以，与其把每个词内部重复录音保得很高，不如把每词配额降一点，再把预算留给更多词种。

\section{如果坚持 8:1:1，怎样分配更合理}

表~\ref{tab:ratio-options} 比较了几种严格 8:1:1 的每词配额。这里“候选词数”指满足该门槛、能进入候选池的英文词种数；“600 词估算体积”用当前子集平均每个 WAV 的大小近似估算。

\begin{table}[H]
    \centering
    \caption{几种严格 8:1:1 配额方案的比较}
    \label{tab:ratio-options}
    \small
    \begin{tabular}{>{\raggedright\arraybackslash}p{2.1cm}cccc}
        \toprule
        每词配额 & 每词总文件数 & 候选词数 & 若选 600 词的总文件数 & 600 词估算体积 \\
        \midrule
        16/2/2 & 20 & 16,111 & 12,000 & 0.36 GiB \\
        24/3/3 & 30 & 13,543 & 18,000 & 0.54 GiB \\
        40/5/5 & 50 & 10,417 & 30,000 & 0.90 GiB \\
        48/6/6 & 60 & 9,395 & 36,000 & 1.08 GiB \\
        64/8/8 & 80 & 7,884 & 48,000 & 1.43 GiB \\
        80/10/10 & 100 & 6,723 & 60,000 & 1.79 GiB \\
        \bottomrule
    \end{tabular}
\end{table}

这张表里最重要的结论是：\textbf{40/5/5 是最均衡的中间点。}

原因如下：
\begin{itemize}[nosep]
    \item 它是严格的 8:1:1；
    \item 每词共 50 条音频，已经足够做训练、调阈值和最终测试；
    \item 候选词仍有 10,417 个，不会把词表压得太狠；
    \item 如果做 600 词子集，总量约 30,000 条、约 0.9 GiB，课程项目可控。
\end{itemize}

相反：
\begin{itemize}[nosep]
    \item \texttt{16/2/2} 或 \texttt{24/3/3} 太薄，dev/test 容易波动；
    \item \texttt{80/10/10} 太重，磁盘和整理成本上去很快，而且额外预算更多花在“同一单词重复录音”，不如拿去换更多词。
\end{itemize}

\section{我们到底应该取多少词}

如果把每词配额固定为 \textbf{40/5/5}，只改变总词数，则不同目标规模的含义如表~\ref{tab:target-scales}。

\begin{table}[H]
    \centering
    \caption{在 40/5/5 下，不同目标词数的含义}
    \label{tab:target-scales}
    \begin{tabular}{ccccp{5.0cm}}
        \toprule
        目标词数 & 总文件数 & 估算体积 & 代表性覆盖 & 适用场景 \\
        \midrule
        160 & 8,000 & 0.24 GiB & 太少 & 不建议，词种仍偏少 \\
        300 & 15,000 & 0.45 GiB & 中等 & 只做四个浊辅音时可接受 \\
        400 & 20,000 & 0.60 GiB & 中等偏稳 & 课程交付版比较合适 \\
        600 & 30,000 & 0.90 GiB & 较稳 & 我最推荐的主实验规模 \\
        800 & 40,000 & 1.19 GiB & 更全面 & 想覆盖更多浊辅音时可考虑 \\
        1000 & 50,000 & 1.49 GiB & 最全面 & 成本更高，课程项目未必需要 \\
        \bottomrule
    \end{tabular}
\end{table}

\subsection{为什么我推荐 600 词，而不是继续用 160 词}

用当前脚本默认的 greedy 选词逻辑，在 \texttt{40/5/5} 下做模拟，几个代表性尾部标签的词数变化如表~\ref{tab:coverage-sim}。

\begin{table}[H]
    \centering
    \caption{40/5/5 下，不同目标词数的尾部标签覆盖模拟}
    \label{tab:coverage-sim}
    \begin{tabular}{ccccccc}
        \toprule
        目标词数 & b & g & jh & y & zh & 说明 \\
        \midrule
        160 & 16 & 11 & 6 & 10 & 4 & 与当前量级接近，仍然偏小 \\
        300 & 31 & 25 & 10 & 17 & 4 & 四标签小实验可用 \\
        400 & 46 & 29 & 16 & 20 & 4 & 比 160 明显稳一些 \\
        600 & 65 & 36 & 25 & 29 & 6 & 绝大多数标签已明显改善 \\
        800 & 82 & 51 & 34 & 39 & 8 & 覆盖进一步提升 \\
        \bottomrule
    \end{tabular}
\end{table}

这说明：
\begin{itemize}[nosep]
    \item 把词数从 160 提到 600，绝大多数标签都会显著改善；
    \item 对课程项目来说，600 已经是一个很好的中间点；
    \item 继续增加到 800 或 1000 当然更全面，但收益和成本开始一起上升。
\end{itemize}

\subsection{还要补一个关键点：不能只改词数，还要改选词策略}

当前脚本的 \path{src/prepare_mswc_subset.py} 在选词时使用
\begin{quote}
\texttt{select\_words(candidates, target\_words, seed\_words\_per\_label=4)}
\end{quote}

也就是说，每个标签最开始只保底拿 4 个词，然后再做全局 greedy fill。这个策略对头部标签影响不大，但会持续压低尾部标签，尤其是 \texttt{/zh/}。

我本地做过一个简单模拟：在 \texttt{40/5/5}、\texttt{600 词} 的设定下：
\begin{itemize}[nosep]
    \item 如果继续用当前 \texttt{seed\_words\_per\_label = 4}，则 \texttt{/zh/} 只有 \textbf{6} 个词；
    \item 如果把每标签保底词数提到 \textbf{20}，则 \texttt{/zh/} 可以提升到 \textbf{20} 个词。
\end{itemize}

所以真正合理的完整方案不是只改成 8:1:1，而是三件事一起改：
\begin{enumerate}[nosep]
    \item 每词比例改成严格 8:1:1；
    \item 总词数从 160 提升到 600 左右；
    \item 把每标签保底词数从 4 提高到更合理的数，例如 20，或者显式加入 \texttt{min\_words\_per\_label} 逻辑。
\end{enumerate}

\section{最终建议}

把上面的分析压缩成一句话：

\begin{quote}
\textbf{MSWC English 原始数据完全够大；当前项目的问题不是数据源太小，而是子集只取了 160 个词。若要改成 8:1:1，我建议优先采用 40/5/5，并把词数提升到 600；如果还想覆盖更多浊辅音，就同步提高每个标签的保底词数。}
\end{quote}

更具体一点：
\begin{table}[H]
    \centering
    \caption{按任务目标给出的推荐规模}
    \label{tab:recommend}
    \begin{tabularx}{\textwidth}{>{\raggedright\arraybackslash}p{4.6cm}cc>{\raggedright\arraybackslash}X}
        \toprule
        目标 & 推荐词数 & 推荐配额 & 说明 \\
        \midrule
        只做四个浊辅音的课程主线 & 300--400 & 40/5/5 & 比当前 160 词稳很多，适合先完成交付 \\
        想扩到更多浊辅音 & \textbf{600} & \textbf{40/5/5} & 我认为这是最均衡的主方案 \\
        想尽量覆盖全部可用浊辅音并做更稳的 AI 比较 & 800--1000 & 40/5/5 或 48/6/6 & 更全面，但成本更高 \\
        \bottomrule
    \end{tabularx}
\end{table}

因此，本文给出的主结论是：
\begin{quote}
\textbf{推荐主方案：600 词 + 40/5/5 + 提高每标签保底词数。}
\end{quote}

\section*{附：建议命令}

如果后续要落地到新的子集，我建议先以如下命令为目标配置：

\begin{verbatim}
cd 216/project/topic2_speech
.venv/bin/python src/prepare_mswc_subset.py \
  --out-root data/raw/mswc_en_subset_600 \
  --target-words 600 \
  --train-per-word 40 \
  --dev-per-word 5 \
  --test-per-word 5 \
  --seed-words-per-label 20
\end{verbatim}

上面这一步现在已经完成：当前仓库里的 \path{src/prepare_mswc_subset.py} 默认就会生成 600 词版本，并且支持显式传入 \texttt{--seed-words-per-label 20}。

\end{document}
